\documentclass[serif, aspectratio=169]{beamer}
%\documentclass[serif]{beamer}  % for 4:3 ratio
\usepackage[T1]{fontenc} 

%\renewcommand{\familydefault}{\sfdefault} % make body text use sans globally
\usepackage{kotex}
%\usepackage{fontspec}
%\setsansfont[
%Path = fonts/,
%UprightFont = *Medium.ttf,
%BoldFont    = *Bold.ttf,
%AutoFakeSlant = 0.2
%]{GmarketSansTTF}
%\setmainhangulfont[
%Path = fonts/,
%UprightFont = *Medium.ttf,
%BoldFont    = *Bold.ttf,
%AutoFakeSlant = 0.2
%]{GmarketSansTTF}

\renewcommand*\familydefault{\sfdefault} %% Only if the base font of the document is to be sans serif
\usepackage[T1]{fontenc}
\usepackage[utf8]{inputenc}
\usepackage{avant}

%\usepackage{fourier} % see "http://faq.ktug.org/wiki/uploads/MathFonts.pdf" for other options
\usepackage{hyperref}
\usepackage{latexsym,amsmath,xcolor,multicol,booktabs,calligra}
\usepackage{graphicx,pstricks,listings,stackengine}
\usepackage{lipsum}
\usepackage{soul}
\usepackage{bm}

\usepackage{tikz}
\usetikzlibrary{arrows.meta, positioning, calc}

\author{Instructor: \underline{Ji Yonghyeon}\\ [4pt]
	\small\texttt{hacker3740@kookmin.ac.kr} }
%\title{Analysis of libOQS:}
\title[The PQC Hard Problem Colloquium]{\textbf{The PQC Hard Problem Colloquium}}
\subtitle{Lattice, Code, Hash, Multivariate, Isogeny, MPC/ZK}
\institute{
	Department of Cyber Security \\
	Koomin Universeity \\
}
%\date{\small \today}
\date{\today}
\usepackage{HKUSTstyle}

% defs
\def\cmd#1{\texttt{\color{red}\footnotesize $\backslash$#1}}
\def\env#1{\texttt{\color{blue}\footnotesize #1}}
% set colors
\definecolor{hkustyellow}{RGB}{167, 131, 55}
\definecolor{hkustblue}{RGB}{0, 56, 116}
\definecolor{hkustred}{RGB}{209, 51, 59}


%\lstset{
%	basicstyle=\ttfamily\small,
%	keywordstyle=\bfseries\color{deepblue},
%	emphstyle=\ttfamily\color{deepred},    % Custom highlighting style
%	stringstyle=\color{deepgreen},
%	numbers=left,
%	numberstyle=\small\color{halfgray},
%	rulesepcolor=\color{red!20!green!20!blue!20},
%	frame=shadowbox,
%}
% ====== LISTINGS (code style) ======
\lstset{
	basicstyle=\ttfamily\small,
	breaklines=true,
	frame=single,
%	backgroundcolor=\color{gray!5},
	tabsize=4,
	keywordstyle=\color{blue!70!black},
	commentstyle=\color{green!40!black},
	columns=fullflexible,
	keepspaces=true,
	upquote=true,
	showstringspaces=false,
	numbers=left,
	numberstyle=\tiny, numbersep=6pt
}


%- --- --- --- --- --- --- --- --- --- --- --- --- --- --- --- 
\begin{document}
	
	\begin{frame}
		\titlepage
		\vspace*{-0.6cm}
		\begin{figure}[htpb]
			\begin{center}
				\includegraphics[keepaspectratio, scale=0.02]{pic/UST}
			\end{center}
		\end{figure}
	\end{frame}
	
	\begin{frame}    
		\tableofcontents[sectionstyle=show,
		subsectionstyle=show/shaded/hide,
		subsubsectionstyle=show/shaded/hide]
	\end{frame}

\section{Introduction}
	\begin{frame}{Purpose of This Colloquium}
	\begin{itemize}
		\item This colloquium is about the \textbf{hard problems} that underlie modern public-key cryptography and post-quantum cryptography.
		\item[]
		\item The goal is not only to know \emph{which schemes exist}, but to understand:
		\begin{itemize}
			\item what mathematical problem each scheme relies on,
			\item what the best known attacks are,
			\item how classical and quantum complexity differ,
%			\item and why some families remain viable while others have failed.
		\end{itemize}
%		\item In other words, the central object of study is the \textbf{security assumption itself}.
	\end{itemize}
\end{frame}

\section{Classical Foundations}

\begin{frame}{Classical Number-Theoretic Foundations}
	Historically, public-key cryptography was built on three dominant assumptions:
	\begin{itemize}
		\item \textbf{Integer Factorization} \hfill RSA, Rabin
		\item \textbf{Discrete Logarithm in Finite Fields} \hfill Diffie--Hellman, DSA, ElGamal
		\item \textbf{Elliptic Curve Discrete Logarithm} \hfill ECDH, ECDSA, EdDSA
	\end{itemize}
	
%	\vspace{0.5em}
%	These problems were attractive because they supported efficient and fast public-key constructions.
\end{frame}

	\begin{frame}{Integer Factorization}
	\begin{itemize}
		\item Integer factorization underlies \textbf{RSA} and \textbf{Rabin}. \item[]
		\item Classically, the best known attacks include:
		\begin{itemize}
			\item \textbf{Quadratic Sieve},
			\item \textbf{General Number Field Sieve}.
		\end{itemize} \item[]
%		\item These are \textbf{sub-exponential} classical attacks. \item[]
		\item In the quantum setting, \textbf{Shor's algorithm} solves the problem in polynomial time.
	\end{itemize}
	
%	\begin{block}{Orientation Point}
%		Factorization is a model example of a problem that is hard enough classically, yet fundamentally broken quantumly.
%	\end{block}
\end{frame}

\begin{frame}{Discrete Logarithm and Elliptic Curve DLP}
	\textbf{Finite-field DLP}
	\begin{itemize}
		\item Supports Diffie--Hellman, DSA, and ElGamal.
		\item Best classical attacks include \textbf{Index Calculus} and \textbf{Number Field Sieve}.
%		\item These are again sub-exponential classically.
	\end{itemize}
	
	\vspace{0.5em}
	\textbf{Elliptic-curve DLP}
	\begin{itemize}
		\item Supports ECDH, ECDSA, and EdDSA.
		\item Generic classical attacks are roughly \textbf{$O(\sqrt{n})$}, e.g.\ Pollard's rho.
		\item Yet both finite-field DLP and ECDLP fall to \textbf{Shor's algorithm} in polynomial time.
	\end{itemize}
\end{frame}

\begin{frame}{The Transition to PQC}
	\begin{block}{Key Transition}
		The post-quantum transition is the replacement of Shor-vulnerable assumptions with alternative assumptions for which no efficient quantum algorithms are known.
	\end{block}
	
	\begin{itemize}
%		\item This is not merely a standards transition.
		\item It is a transition in \textbf{mathematical foundations}.
		\item The main candidate families are:
		\begin{itemize}
			\item lattices,
			\item codes,
			\item hash-based constructions,
			\item multivariate systems,
			\item isogenies.
		\end{itemize}
	\end{itemize}
\end{frame}

\section{Post-Quantum Hard Problem Families}

\begin{frame}{Lattice Problems}
	The lattice family includes:
	\begin{itemize}
		\item \textbf{LWE},
		\item \textbf{SIS},
		\item \textbf{SVP},
		\item \textbf{CVP}.
	\end{itemize}
	
	Representative systems include:
	\begin{itemize}
		\item Kyber / ML-KEM,
		\item Dilithium / ML-DSA,
		\item NTRU,
		\item fully homomorphic encryption constructions.
	\end{itemize}
\end{frame}

%\begin{frame}{How Lattice Systems Are Attacked}
%	The dominant attack language for lattice-based cryptography includes:
%	\begin{itemize}
%		\item \textbf{Lattice reduction}, especially LLL and BKZ,
%		\item \textbf{Sieve algorithms},
%		\item careful \textbf{concrete security estimation}.
%	\end{itemize}
	
%	\vspace{0.5em}
%	The best known attacks remain broadly \textbf{exponential} in the relevant dimension parameter.
%	
%	\begin{block}{Orientation Point}
%		Lattice cryptography is currently the most important practical PQC family precisely because no Shor-style collapse is known.
%	\end{block}
%\end{frame}

\begin{frame}{Code-Based Problems}
	Code-based cryptography is centered on:
	\begin{itemize}
		\item \textbf{syndrome decoding},
		\item \textbf{quasi-cyclic syndrome decoding problem}
		\item more generally, decoding random linear codes.
	\end{itemize}
	
	Representative systems include:
	\begin{itemize}
		\item McEliece,
		\item HQC,
		\item BIKE.
	\end{itemize}
	
%	The dominant attack family is \textbf{Information Set Decoding (ISD)} and its variants.
\end{frame}

%\begin{frame}{Code-Based Security Perspective}
%	\begin{itemize}
%		\item Best known attacks remain \textbf{exponential}.
%		\item Quantum improvements are comparatively limited.
%		\item Grover's algorithm gives only a \textbf{quadratic speedup}, rather than a full structural break.
%	\end{itemize}
%	
%	\begin{block}{Orientation Point}
%		Code-based cryptography remains one of the most conservative and durable post-quantum directions.
%	\end{block}
%\end{frame}
%
\begin{frame}{Hash-Based Assumptions}
	Hash-based cryptography relies on:
	\begin{itemize}
		\item collision resistance,
		\item preimage resistance,
		\item second-preimage resistance.
	\end{itemize}
	
	Representative signature systems include:
	\begin{itemize}
		\item SPHINCS+,
		\item XMSS,
		\item LMS.
	\end{itemize}
\end{frame}
%
%\begin{frame}{Hash-Based Security Perspective}
%	\begin{itemize}
%		\item Collision attacks follow the \textbf{birthday bound}.
%		\item Quantum preimage search is accelerated by \textbf{Grover's algorithm}.
%		\item But there is no analogue of Shor's algorithm causing polynomial-time collapse.
%	\end{itemize}
%	
%	\begin{block}{Orientation Point}
%		Hash-based cryptography is often viewed as assumption-conservative, even when it is less attractive in size or efficiency.
%	\end{block}
%\end{frame}
%
\begin{frame}{Multivariate Problems}
	Multivariate cryptography is based on the difficulty of solving systems of quadratic equations over finite fields.
	
	\vspace{0.5em}
	This is commonly framed as the \textbf{MQ problem}.
	
	\vspace{0.5em}
	Representative systems include:
	\begin{itemize}
		\item UOV,
		\item MAYO,
		\item historically, Rainbow.
	\end{itemize}
\end{frame}
%
%\begin{frame}{Multivariate Attack Perspective}
%	Relevant attack language includes:
%	\begin{itemize}
%		\item \textbf{MinRank},
%		\item \textbf{Kipnis--Shamir},
%		\item \textbf{XL / Gr\"obner basis} methods,
%		\item and scheme-specific structural attacks.
%	\end{itemize}
%	
%	\vspace{0.5em}
%	While the general problem remains hard-looking, particular constructions may fail because they reveal exploitable algebraic structure.
%\end{frame}
%
%\begin{frame}{The Rainbow Lesson}
%	\begin{itemize}
%		\item Rainbow was a prominent multivariate signature construction.
%		\item It was broken by a polynomial-time structural attack.
%		\item This teaches an important lesson:
%	\end{itemize}
%	
%	\begin{block}{Critical Distinction}
%		Family-level hardness does not automatically imply scheme-level security.
%	\end{block}
%\end{frame}
%
\begin{frame}{Isogeny Problems}
	Isogeny-based cryptography is among the most mathematically elegant areas in PQC.
	
	\vspace{0.5em}
	Representative names include:
	\begin{itemize}
		\item SQISign,
		\item CSIDH,
		\item historically, SIDH / SIKE.
	\end{itemize}
\end{frame}
%
%\begin{frame}{The SIKE Lesson}
%	\begin{itemize}
%		\item SIKE was broken by the \textbf{Castryck--Decru attack}.
%		\item This attack was polynomial time and completely changed the status of SIDH/SIKE.
%		\item Other attack frameworks include meet-in-the-middle methods and quantum algorithms such as Kuperberg's algorithm in certain commutative settings.
%	\end{itemize}
%	
%	\begin{block}{Orientation Point}
%		Isogeny cryptography remains mathematically rich, but it must now be studied with exceptional caution regarding hidden structure.
%	\end{block}
%\end{frame}
%
%\section{How to Read the Hard-Problem Landscape}
%
%\begin{frame}{Comparative Reading of the Landscape}
%	The families divide naturally into four security narratives:
%	\begin{itemize}
%		\item \textbf{Factorization / DLP / ECDLP}: classically strong, quantumly broken.
%		\item \textbf{Lattice / Code-based}: currently the strongest practical PQC foundations.
%		\item \textbf{Hash-based}: conservative and structurally simple.
%		\item \textbf{Multivariate / Isogeny}: mathematically compelling, but more exposed to structural fragility.
%	\end{itemize}
%\end{frame}
%
%\begin{frame}{Three Questions for Every Scheme}
%	For every construction discussed in this colloquium, we should ask:
%	\begin{enumerate}
%		\item What hard problem is assumed?
%		\item What are the best known classical and quantum attacks?
%		\item Does the scheme introduce exploitable structure beyond the underlying hard problem?
%	\end{enumerate}
%	
%	\vspace{0.75em}
%	This provides a disciplined framework for comparing PQC proposals.
%\end{frame}

%\begin{frame}{What the Hard-Problem Summary Really Teaches}
%	The summary should not be read merely as a catalog of attacks.
%	
%	\vspace{0.5em}
%	It should be read as a set of research principles:
%	\begin{itemize}
%		\item distinguish \textbf{problem hardness} from \textbf{scheme security},
%		\item distinguish \textbf{generic attacks} from \textbf{structural attacks},
%		\item compare \textbf{classical} and \textbf{quantum} attack models,
%		\item and study broken designs as lessons in cryptographic methodology.
%	\end{itemize}
%\end{frame}
%%
\section{The Colloquium}
%
\begin{frame}{Suggested Colloquium Flow}
	A natural progression for the colloquium is:
	\begin{enumerate}
		\item Classical number-theoretic assumptions and Shor's algorithm
		\item Lattice problems and reduction-based security
		\item Code-based assumptions and ISD
		\item Hash-based signatures %as a conservative baseline
		\item Multivariate constructions %and structural failure
		\item Isogeny cryptography %after SIKE
		\item Cross-family comparison of security assumptions
	\end{enumerate}
\end{frame}

\begin{frame}{Expected Outcomes}
	By the end of the colloquium, participants should be able to:
	\begin{itemize}
		\item[]
		\item explain the main hard-problem families in PQC,
		\item[]
		\item describe the best known attacks against each family.
%		\item understand why some assumptions remain credible,
%		\item and evaluate a cryptosystem at the level of its actual security assumption.
	\end{itemize}
\end{frame}

\begin{frame}{Closing}
%	\begin{block}{}
%		Post-quantum cryptography is not simply the search for new algorithms. It is the search for mathematical problems whose hardness survives quantum computation, structural cryptanalysis, and concrete attack advances.
%	\end{block}
		\begin{itemize}
%		\item[]
		\item Post-quantum cryptography is not simply the search for new algorithms.
		\item[]
		\item It is the search for mathematical problems whose hardness survives quantum computation, structural cryptanalysis, and attack advances.
	\end{itemize}

	\vspace{0.75em}
	That is the central purpose of the PQC Hard Problem Colloquium.
\end{frame}

\section{Preview}

\begin{frame}{A General Cryptographic Template}
	Let
	\begin{itemize}
		\item $\mathcal{I}$ is the public instance,
		\item $\mathcal{S}$ is the hidden witness, secret, or structure,
%		\item and the adversary is asked to recover $\mathcal{S}$, or to solve a decisional variant.
	\end{itemize}
	
	\vspace{0.5em}
	A hard problem is typically one of the following:
	\[
	\text{Search: } \mathcal{A}(\mathcal{I}) = \mathcal{S},
	\qquad
	\text{Decision: } \mathcal{A}(\mathcal{I}) \in \{0,1\},
	\]
	with no efficient known algorithm for the relevant parameter regime.
\end{frame}

\begin{frame}{Integer Factorization}
	\textbf{Instance.}
	Let
	\[
	N = pq,
	\]
	where $p,q$ are large primes. The public input is $N$.
	
	\vspace{2.4em}
	\textbf{Search problem.}
	Recover nontrivial factors:
	\[
	\text{Given } N, \text{ find } p,q \text{ such that } N=pq.
	\]
\end{frame}

\begin{frame}{Discrete Logarithm Problem in Finite Fields}
	Let $G = \langle g \rangle$ be a cyclic group, typically
	\[
	G = \mathbb{F}_p^\times.
	\]
	
	\textbf{Instance.}
	Given
	\[
	g \in G, \qquad h = g^x \in G.
	\]
	
	\vspace{2.4em}
	\textbf{Search problem.}
	Recover the exponent $x$:\[
	\text{find } x \text{ such that } h=g^x.
	\]
\end{frame}

\begin{frame}{Elliptic Curve Discrete Logarithm Problem (ECDLP)}
	Let $E/\mathbb{F}_q$ be an elliptic curve, and let $
	P \in E(\mathbb{F}_q)$
	be a point of large order $n$.
	
	\vspace{2.4em}
	\textbf{Instance.}
	Given
	\[
	P, \qquad Q = xP.
	\]
	
	\vspace{2.4em}
	\textbf{Search problem.}
	Recover
	\[
	x \in \mathbb{Z}_n
	\quad \text{such that} \quad
	Q = xP.
	\]
\end{frame}

	
\begin{frame}{Lattice Problems}
	A lattice generated by basis vectors $\bm{b}_1,\dots,\bm{b}_m \in \mathbb{R}^n$ is
	\[
	\Lambda = \left\{ \sum_{i=1}^m z_i \bm{b}_i : z_i \in \mathbb{Z} \right\}.
	\]
	
	\textbf{Shortest Vector Problem (SVP).}
	Find
	\[
	\bm{v} \in \Lambda \setminus \{\bm{0}\}
	\]
	minimizing
	\[
	\|\bm{v}\|.
	\]
	
	\textbf{Closest Vector Problem (CVP).}
	Given $\bm{t} \in \mathbb{R}^n$, find
	\[
	\bm{v} \in \Lambda
	\]
	minimizing
	\[
	\|\bm{t}-\bm{v}\|.
	\]
\end{frame}


	\begin{frame}{Learning With Errors (LWE)}
%	LWE is an average-case algebraic-noise problem.
	
	\textbf{Secret.}
	$
	\bm{s} \in \mathbb{Z}_q^n.
	$
	
	\vspace{0.4em}
	\textbf{Samples.}
	For random
	$
	\bm{a}_i \in \mathbb{Z}_q^n,
	$
	we observe
	\[
	b_i = \langle \bm{a}_i,\bm{s} \rangle + e_i \pmod q,
	\]
	where $e_i$ is a small error.
	
	\vspace{0.4em}
	\textbf{Search-LWE.}
	Given many pairs
	\[
	(\bm{a}_i,b_i),
	\]
	recover $\bm{s}$.
	
	\textbf{Decision-LWE.}
	Distinguish
	\[
	(\bm{a}_i,\langle \bm{a}_i,\bm{s}\rangle + e_i)
	\]
	from uniform samples in
	$
	\mathbb{Z}_q^n \times \mathbb{Z}_q.
	$
	
%	\vspace{0.4em}
%	\textbf{Symbolic intuition.}
%	Without noise, one solves a linear system.
%	With noise, the system becomes
%	\[
%	A\bm{s} + \bm{e} \equiv \bm{b} \pmod q,
%	\]
%	and the small perturbation $\bm{e}$ is exactly what creates hardness.
\end{frame}

%------------------------------------------------

\begin{frame}{SIS: Short Integer Solution}
	Let
	\[
	A \in \mathbb{Z}_q^{m \times n}.
	\]
	
	\textbf{Problem.}
	Find a nonzero vector
	\[
	\bm{x} \in \mathbb{Z}^n
	\]
	such that
	$
	A\bm{x} \equiv \bm{0} \pmod q
	$
	and
	$
	\|\bm{x}\|
	$
	is small.
	
%	\vspace{2.4em}
%	\textbf{Interpretation.}
%	SIS asks for a short relation among the columns of $A$ modulo $q$.
\end{frame}

	\begin{frame}{Code-Based Hardness: Syndrome Decoding}
	Let
	\[
	H \in \mathbb{F}_2^{r \times n}
	\]
	be a parity-check matrix.
	
	\vspace{1.2em}
	\textbf{Given.}
	A syndrome
	\[
	\bm{s} \in \mathbb{F}_2^r
	\]
	and a weight bound $w$.
	
	\vspace{1.2em}
	\textbf{Problem.}
	Find an error vector
	\[
	\bm{e} \in \mathbb{F}_2^n
	\]
	such that $
	H\bm{e}^T = \bm{s}^T$
	and
	\[
	\mathrm{wt}(\bm{e}) \le w.
	\]
	
%	\vspace{0.4em}
%	\textbf{Interpretation.}
%	The vector $\bm{e}$ is sparse, but hidden by the linear map $H$.
%	
%	\vspace{0.4em}
%	\textbf{Cryptographic asymmetry.}
%	Computing
%	\[
%	\bm{s} = H\bm{e}^T
%	\]
%	is easy; recovering a low-weight $\bm{e}$ is hard.
%	
%	\vspace{0.4em}
%	\textbf{Attack language.}
%	\[
%	\text{ISD},\ \text{Stern},\ \text{Lee--Brickell},\ \text{BJMM-type methods}.
%	\]
\end{frame}

\begin{frame}{Hash-Based Hardness}
	Let
	\[
	H : \{0,1\}^\ast \to \{0,1\}^n
	\]
	be a cryptographic hash function.
	
	\vspace{1.2em}
	\textbf{Preimage resistance.}
	Given
	$
	y \in \{0,1\}^n,
	$
	find
	\[
	x \text{ such that } H(x)=y.
	\]
	
	\vspace{1.2em}
	\textbf{Second-preimage resistance.}
	Given $x$, find
	\[
	x' \neq x
	\quad \text{such that} \quad
	H(x') = H(x).
	\]
	
	\vspace{1.2em}
	\textbf{Collision resistance.}
	Find
	\[
	x \neq x'
	\quad \text{such that} \quad
	H(x)=H(x').
	\]
	
%	\vspace{0.4em}
%	\textbf{Generic bounds.}
%	\[
%	\text{collision search: } O(2^{n/2}),
%	\qquad
%	\text{preimage search: } O(2^n).
%	\]
	
%	\textbf{Quantum effect.}
%	Grover reduces exhaustive preimage search to roughly
%	\[
%	O(2^{n/2}).
%	\]
\end{frame}
	
	% --- Thank you slide ---
	\begin{frame}
		\begin{center}
			{ Thank you for listening !}
			\vspace{1cm}
			
			Ji, YongHyeon \\[1em]
			\texttt{hacker3740@kookmin.ac.kr} 
		\end{center}
	\end{frame}
	
\end{document}
